% ===================================================================
%  اللوحة رقم ٤ — سنّ النضج والشيخوخة.
%
%  Traduction, faite en 2026, en arabe levantin d'aujourd'hui, sur
%  l'ido de texto/io. Regles de la colonne : voir 10-lawha-01.tex.
%  Deux scenes, deux numerotations.
%
%  LE COUSIN DU LIVRET NE PEUT PAS S'ECRIRE EN ARABE SANS QU'ON
%  DECIDE DE QUI IL EST LE FILS. « kuzulo » et « kuzino » ne se
%  traduisent pas : il faut choisir entre ابن العمّ, le fils du frere
%  du pere, et ابن الخال, le fils du frere de la mere — deux mots que
%  rien ne recouvre, dans une langue ou l'oncle paternel et l'oncle
%  maternel ne portent pas le meme nom depuis toujours. LE RELEVE
%  TRANCHE UNE FOIS POUR TOUTES ET DU COTE DU PERE : c'est le defaut
%  d'usage quand rien n'est dit, et le tableau 3 avait deja ecrit
%  « ابن العمّ » pour la meme famille. On ne peut pas rendre au
%  lecteur l'ignorance ou le francais le laisse ; on peut au moins ne
%  pas changer d'avis d'un tableau a l'autre.
%
%  ET LA OU LE FRANCAIS DIT « BEAU- » CINQ FOIS, LE LEVANTIN DIT CINQ
%  MOTS DIFFERENTS. Le beau-pere (20) est « الحمو », la belle-mere
%  (21) « الحماة », la belle-fille « كنّة », la belle-sœur « سلفة », le
%  beau-frere (30) « النسيب ». Aucun n'est derive d'un autre, aucun ne
%  se devine : ce sont cinq places dans une maison, et la langue les
%  a nommees separement parce qu'on y vit separement.
%
%  LE GARCON D'HONNEUR (25) EST « إشبين », ET C'EST DU SYRIAQUE. Le
%  cure du tableau 3 etait « الخوري », syriaque lui aussi. LE
%  VOCABULAIRE CHRETIEN DU LEVANT EST DU SYRIAQUE SOUS DE L'ARABE, et
%  le livret, qui grave un bapteme de village puis une noce de
%  village, tombe deux fois de suite exactement dessus.
%
%  LES MOIS CONTINUENT CE QUE LE TABLEAU 3 AVAIT COMMENCE : octobre
%  est « تشرين الأول », novembre « تشرين التاني », decembre « كانون
%  الأول ». Mais le mot qui les contient (43) est « روزنامة », qui est
%  persan — rūznāma, le livre des jours. LES MOIS SONT AKKADIENS, LE
%  MOT QUI LES TIENT EST PERSAN, ET LA LANGUE EST ARABE.
%
%  Enfin la garde-robe de ce tableau est un musee. Le pardessus (7)
%  est « بالطو » et le fauteuil (34) « فوتيه », tous les deux
%  francais ; le gateau (46) est « كاتو » ; la ceinture (5) est
%  « زنّار », qui est grec. Et la neige (19) est « التلج » et non
%  الثلج : le ث tombe dans le ت pour la quatrieme fois depuis le
%  tableau 2, apres le dos, le lundi et le taureau.
% ===================================================================

%%K t04-apar-1 apar
\VUsaut{4.0mm}
\VUtitre{15.0pt}{60}{اللوحة رقم ٤}
\VUsaut{1.6mm}
\VUtitre{11.4pt}{0}{سنّ النضج والشيخوخة.}
\VUsaut{3.2mm}

%%K t04-tit-1 sub
\VUcentre{11.4pt}{0}{\textit{المشهد الأول.}}
\VUsaut{1.6mm}
\VUtitre{11.4pt}{0}{عزيمة العرس.}
\VUsaut{2.6mm}

%%K t04-tit-2 sub
\VUpk{10.2pt}{40}{الخريف.}
\VUsaut{3.0mm}

%%K t04-c1-01-1 p
1. --- الشباب كبروا. وواحد منهن، حنّا، عم يتجوّز. نحنا حاضرين
\VUgras{عزيمة العرس} اللي بتجمع، حوالين نفس الطاولة، عيلتين العروسين.

%%K t04-c1-02-1 p
2. --- نحنا بـ\VUgras{الخريف}، \VUgras{بشهر أيلول}. هوا
\VUgras{تشرين الأول} و\VUgras{تشرين التاني} البارد لسّا ما عرّى الضيعة
من ورقها الأخضر. هوا المسا دافي ومعطّر ؛ لهيك العشا عم يصير تحت
\VUgras{الفرندة} \textsuperscript{(8)} اللي عالجنينة.

%%K t04-c1-03-1 p
3. --- وبعيد، عالـ\VUgras{تلّة} \textsuperscript{(3)}، بيت أهل حنّا،
\VUgras{قصر} \textsuperscript{(1)} أنيق مخبّى بالخضرة ؛ وكروم عنب حلوة،
بتعطي نبيد ممتاز، مغطّية منحدر التلّة. وصاحب الكرم عم يفلحها ويعتني
فيها.

%%K t04-c1-04-1 p
4. --- وفوق الكروم، عم يمرق سرب \VUgras{حجل} \textsuperscript{(2)}
مذعور من اقتراب \VUgras{حارس الصيد} \textsuperscript{(5)}. وهاد، مع
\VUgras{بارودته} تحت باطه و\VUgras{كيس الصيد} معلّق عكتفه، عم يفتّش
\VUgras{الأحراش} \textsuperscript{(4)} مع \VUgras{كلبه} \textsuperscript{(6)}.
عم يراقب الملك وبيمنع الصيّادين المخالفين يقضوا عالصيد.

%%K t04-c1-05-1 p
5. --- وبرّات \VUgras{الفرندة}، عم تتسلّق \VUgras{تعريشة عنب} معلّقة
منها \VUgras{عناقيد} \textsuperscript{(7)} كتيفة.

%%K t04-c1-06-1 p
6. --- وزهور الخريف الحلوة اللي مزيّنة \VUgras{الطاولة}
\textsuperscript{(16)} هيّي \VUgras{الأقحوان} \textsuperscript{(9)} ؛
و\VUgras{أبو} \textsuperscript{(10)} العروس حطّ وحدة منهن بعروة
\VUgras{فراكه}.

%%K t04-c1-07-1 p
7. --- الأكل عم يخلص. \VUgras{الخادم} \textsuperscript{(11)} بلّش يجيب
صحون الحلو وعن قريب رح يشيل قطع الأدوات اللي ما عاد إلها لزوم :
\VUgras{ملاعق} \textsuperscript{(14)}، و\VUgras{شوك} \textsuperscript{(12)}،
و\VUgras{سكاكين} \textsuperscript{(13)}، و\VUgras{صحون كبيرة}
\textsuperscript{(15)}.

%%K t04-tit-3 sub
\VUpk{10.2pt}{40}{الأهل.}
\VUsaut{3.0mm}

%%K t04-c1-08-1 p
8. --- الكل مبيّنين مبسوطين، والابتسامة اللي فيها \VUgras{إمّ}
\textsuperscript{(17)} العريس عم تتطلّع على ولادها بتدلّ على سعادتها.

%%K t04-c1-09-1 p
9. --- هالجوازة بتجمع عيلتين غنيّات بالبلد. حنّا، \VUgras{العريس}
\textsuperscript{(18)}، \VUgras{ابن} صاحب ملك كبير، وعم يصير
\VUgras{صهر} صناعي شاطر.

%%K t04-c1-10-1 p
10. --- \VUgras{العروسين} شباب، ومع هيك من زمان كانوا
\VUgras{مخطوبين}. و\VUgras{العروس} \textsuperscript{(19)}، مارتا، كتير
حلوة بـ\VUgras{فستانها الأبيض} المزيّن بـ\VUgras{زهر الليمون}.

%%K t04-c1-11-1 p
11. --- و\VUgras{حماها} \textsuperscript{(20)} كتير مبسوط إنّه صار إلو
\VUgras{كنّة} بهاللطف، و\VUgras{حماة} \textsuperscript{(21)} حنّا عم تبتسم
وهيّي عم تشوف \VUgras{بنتها} هالقدّ حلوة.

%%K t04-c1-12-1 p
12. --- وبطرف الطاولة، قاعدين باقي \VUgras{المعازيم}
\textsuperscript{(22)}، \VUgras{قرايب ورفقات وولاد عمّ وبنات عمّ}.
و\VUgras{رئيس الخدم} \textsuperscript{(23)}، وفوطة تحت باطه، عم يخدمهن
بنشاط.

%%K t04-tit-4 sub
\VUpk{10.2pt}{40}{الرفقات.}
\VUsaut{3.0mm}

%%K t04-c1-13-1 p
13. --- وعيمين الطاولة، هيّي \VUgras{آنسة} \textsuperscript{(24)}،
\VUgras{صاحبة} العروس الشابّة، وبعدين \VUgras{أخوها}، اللي هوّي
\VUgras{إشبينها} \textsuperscript{(25)}. عم يحكي مع \VUgras{إشبينته}
\textsuperscript{(26)}، أخت العريس، اللي بهالجوازة عم تصير
\VUgras{سلفة} مارتا. شابّة حلوة كتير. عليها \VUgras{مصاغ} حلو،
و\VUgras{عقد لولو}، و\VUgras{إسوارة} دهب.

%%K t04-c1-14-1 p
14. --- وجنبهن، السيّد اللي واقف ورافع كاسه هوّي \VUgras{صاحب}
\textsuperscript{(27)} العيلة. وهوّي واحد من \VUgras{شهود العريس}. عم
\VUgras{يشرب نخب} \VUgras{صحّة} العروسين الشباب وبيتمنّالهن سعادة
طويلة. لابس بدلة الحفلات، فراك مع \VUgras{صبابيط ملمّعة}
\textsuperscript{(28)}. وهوّي مرافق \VUgras{الستّ} \textsuperscript{(29)}
اللي هيّي \VUgras{أرملة}.

%%K t04-c1-15-1 p
15. --- والشاب اللي عليه \VUgras{فراك} \textsuperscript{(31)}، مع شوارب
رفيعة سودا، هوّي \VUgras{نسيب} \textsuperscript{(30)} حنّا. مهندس
كبير. وهوّي جنب \VUgras{سيّدة شابّة} \textsuperscript{(33)} كتير أنيقة،
\VUgras{الأخت الكبيرة} لمارتا وإمّ البنت الزغيرة الحلوة اللي جايي،
معطية إيدها لابن عمّها، تقدّم زهرة وتتمنّى إلها \VUgras{مسا الخير}.
هالسيّدة عليها \VUgras{حلقان} كتير حلوة و\VUgras{غطا راس} أنيق.

%%K t04-c1-16-1 p
16. --- وابن العمّ الزغير، اللي شكله غريب بـ\VUgras{بلوزته}
\textsuperscript{(36)} و\VUgras{كمّه} \textsuperscript{(37)} المنفوخ،
كتير مسكين. هوّي \VUgras{يتيم} \textsuperscript{(35)}. وهوّي زغير كتير
فقد أبوه وإمّه. وعمّه هوّي \VUgras{الوصي} \textsuperscript{(38)} عليه،
وضلّ \VUgras{عزابي} ليربّي هالولد المسكين. رجّال طيّب. شوي
\VUgras{أصلع} وما بيشوف منيح ؛ بيستعمل \VUgras{نظّارات}.

%%K t04-tit-5 sub
\VUsaut{2.4mm}
\VUpk{10.2pt}{40}{الحلو.}
\VUsaut{3.0mm}

%%K t04-c1-17-1 p
17. --- ووراء المعازيم، على طاولة مغطّاة بـ\VUgras{شرشف}، محضّر
\VUgras{الحلو} \textsuperscript{(39)}. بطاسة كبيرة، هيّي الفواكه :
\VUgras{دراق} \textsuperscript{(40)}، و\VUgras{إجّاص}
\textsuperscript{(41)}، و\VUgras{تفّاح} \textsuperscript{(42)}،
و\VUgras{مشمش} \textsuperscript{(43)}، و\VUgras{عنب}
\textsuperscript{(44)}. وجنبها، \VUgras{أناناس}
\textsuperscript{(45)}، و\VUgras{كاتو} \textsuperscript{(46)} حلو،
و\VUgras{قناني مشروب} \textsuperscript{(48)} و\VUgras{شمبانيا}
\textsuperscript{(49)}، وكمان أعمدة \VUgras{صحون} \textsuperscript{(47)}
نضيفة.

%%K t04-c1-18-1 p
18. --- \VUgras{الساقي} \textsuperscript{(50)} عم يفتح \VUgras{قنّينة}
\textsuperscript{(52)} نبيد عتيق بـ\VUgras{فتّاحة الفلّين}
\textsuperscript{(51)}. و\VUgras{الفلّينة} \textsuperscript{(53)} مبيّن
إنّه كتير صعب تنشلح. وهالساقي عنده \VUgras{فوطة} \textsuperscript{(54)}
تحت باطه.

%%K t04-c1-19-1 p
19. --- وبعد العشا، السيّدين رح يروحوا على غرفة التدخين، والسيّدات رح
يروحوا يحضّروا حالهن للرقص.

%%K t04-tit-6 sub
\VUsaut{5.0mm}
\VUcentre{11.4pt}{0}{\textit{المشهد التاني.}}
\VUsaut{1.6mm}
\VUpk{11.4pt}{40}{عيد ميلاد الجدّ.}
\VUsaut{2.6mm}

%%K t04-tit-7 sub
\VUpk{10.2pt}{40}{الزيارة.}
\VUsaut{3.0mm}

%%K t04-c2-01-1 p
1. --- مرقت عشر سنين، وأهل حنّا صاروا ختيارية وكتير بيحبّوا أحفادهن
التلاتة، \VUgras{حفيدين} \textsuperscript{(2)} و\VUgras{حفيدة}
\textsuperscript{(3)}. ولأنه اليوم \VUgras{عيد ميلاد الجدّ}، الولاد
جايين يتمنّولهن عيد ميلاد سعيد.

%%K t04-c2-02-1 p
2. --- وبطرس الزغير لسّا كتير زغير، عمره تلات سنين بس. عم يشوف
\VUgras{البسّ} \textsuperscript{(1)} عم يقرّب. بدّه يمسّد \VUgras{فروته}
الحريرية الحلوة و\VUgras{ديله} الطويل، وما عم يفكّر إنّه تحت
\VUgras{إجريه} الرشيقة مخبّاية ضوافر مدبّبة وخبيثة يمكن تخربشه.

%%K t04-c2-03-1 p
3. --- وأخته غابرييلا، اللي عم تعطيه إيدها، أعقل منه بكتير ؛ ومارسيلوس
مع \VUgras{معطفه} \textsuperscript{(4)} الكبير و\VUgras{زنّاره}
\textsuperscript{(5)} الجلد مبيّن إنّه صار رجّال شوي، مع إنّه كمّه
القصير عم يبيّن \VUgras{كلساته} \textsuperscript{(6)}.

%%K t04-c2-04-1 p
4. --- وأبوهن حطّ \VUgras{بالطوه} \textsuperscript{(7)} التخين تبع
الشتوية مع \VUgras{ياقة الفرو} \textsuperscript{(8)}، وبـ\VUgras{جيبته}
\textsuperscript{(9)} في منديل. ومرته عليها برنيطة اللبّاد المزيّنة
بـ\VUgras{ريش} \textsuperscript{(10)}، و\VUgras{جاكيتها}
\textsuperscript{(11)}، و\VUgras{لفحة الفرو} \textsuperscript{(12)}
و\VUgras{كيس الإيدين} \textsuperscript{(13)}.

%%K t04-c2-05-1 p
5. --- الدني كتير باردة، لأن الشتوية عم تحكم. \VUgras{التلج}
\textsuperscript{(19)} ضلّ نازل طول النهار وسطوح البيوت كلها بيضا. ومع
\VUgras{الليل}، البرد عم يصير أقسى. وبالسما، \VUgras{القمر}
\textsuperscript{(17)} و\VUgras{النجوم} \textsuperscript{(18)} عم يلمعوا
وعم يخترقوا \VUgras{العتمة}.

%%K t04-tit-8 sub
\VUsaut{4.2mm}
\VUpk{10.2pt}{40}{المرض.}
\VUsaut{3.0mm}

%%K t04-c2-06-1 p
6. --- و\VUgras{الأعمى} \textsuperscript{(20)} المسكين اللي عم يقطع
الساحة قدّام \VUgras{الصيدلية} \textsuperscript{(16)} وعم يدلّه
\VUgras{كلبه} \textsuperscript{(21)}، كتير تعيس. ويوستينيا،
\VUgras{الخدّامة} \textsuperscript{(14)}، عم تجيب من المطبخ
\VUgras{زبدية} \textsuperscript{(15)} \VUgras{زهورات} كتير سخنة ومحلّاية
منيح.

%%K t04-c2-07-1 p
7. --- و\VUgras{الستّ} \textsuperscript{(23)}، اللي بدّها تدوّر على
\VUgras{نظّاراتها} \textsuperscript{(26)} عالطاولة لتقرا \VUgras{الوصفة}
\textsuperscript{(27)} اللي كتبها \VUgras{الحكيم} \textsuperscript{(25)}
هلّق، عم تقوم من فوتيهها متّكية على \VUgras{عصاتها}
\textsuperscript{(24)}.

%%K t04-c2-08-1 p
8. --- كتير مرّات بتعاني من \VUgras{أوجاع} و\VUgras{دوخة}
و\VUgras{رشح} و\VUgras{وجع سنان}، وإجريها ضعيفة وعيونها ما عادت
منيحة. ومع هيك، بكل رغبة عم تتمسخر على \VUgras{حكيمها} العتيق اللي
دايماً بدّه يكتبلها \VUgras{شربة دوا} \textsuperscript{(28)}. واليوم
كتير مبسوطة لأن عيلتها الشابّة حواليها.

%%K t04-c2-09-1 p
9. --- والغرفة مسكّرة منيح ودافية. وبـ\VUgras{المدفأة}
\textsuperscript{(29)} الرخام عم تشتعل \VUgras{نار حلوة}
\textsuperscript{(30)}، وبيحسّ الواحد حاله مبسوط بـ\VUgras{ضو}
\VUgras{اللمبة} \textsuperscript{(31)} الحلو اللي هلّق ولّعوها.

%%K t04-tit-9 sub
\VUsaut{4.2mm}
\VUpk{10.2pt}{40}{الجدّ.}
\VUsaut{3.0mm}

%%K t04-c2-10-1 p
10. --- وحتى \VUgras{الجدّ} \textsuperscript{(33)} نسي إنّه كان مريض،
وإنّه \VUgras{الروماتيزم} مثبّته على \VUgras{فوتيهه}
\textsuperscript{(34)}، وإنّه مبارح لسّا كان عنده \VUgras{سخونة
وإغماء}. ما عاد بيتذكّر إنّه بالشتوية الماضية عانى من \VUgras{ذات
الرئة} وإنّه \VUgras{سعال} خشن ومتعب خلّاه يتعذّب كتير.

%%K t04-c2-10-2 p
ومع هيك كتير صعب تعافى. من هلّق صار أحسن شوي. بس إجره اللي عم يسندها
على \VUgras{مخدّة} \textsuperscript{(38)} محطوطة على \VUgras{كرسي زغير}
\textsuperscript{(39)} كمان عم توجعه.

%%K t04-c2-11-1 p
11. --- ولمّا بيكون ملفوف منيح بـ\VUgras{روبه} \textsuperscript{(35)}
الدافي مع \VUgras{زراير} \textsuperscript{(36)} تخينة من صدف، ولمّا
بيكون جنبه، لتقرالو الجريدة، \VUgras{اليتيمة} \textsuperscript{(40)}
الزغيرة لابسة \VUgras{طاقية} بيضا و\VUgras{فستان} \textsuperscript{(42)}
أزرق و\VUgras{عباية} \textsuperscript{(41)}، ما بيشتكي.

%%K t04-c2-12-1 p
12. --- وكل سنة، لمّا \VUgras{الروزنامة} \textsuperscript{(43)} بترجع
تدلّ على \VUgras{تاريخ} أول \VUgras{كانون الأول}، بيستنّى
\VUgras{أحفاده} بفارغ الصبر، لأنه بيعرف إنّهن ما رح ينسوا هالـ\VUgras{تاريخ}
المهمّ.

%%K t04-c2-13-1 p
13. --- وبيتذكّر بتأثّر هديك الأيام البعيدة لمّا كان يروح بنفسه يتمنّى
عيد سعيد لـ\VUgras{المشير} الإمبراطوري المجيد، اللي \VUgras{صورته}
\textsuperscript{(44)} معلّقة عالحيط، واللي هوّي \VUgras{جدّ جدّ}
ولاده.

\VUsaut{6.0mm}
\VUfilet{22mm}
